% Speaker companion for milestone2 — 15-minute presentation.
% Per-slide talking points, background context, and Q&A backup material.
% Compile: pdflatex milestone2_speaker_companion.tex

\documentclass[11pt,a4paper]{article}

\usepackage[margin=2cm,top=1.8cm,bottom=2cm]{geometry}
\usepackage[hidelinks]{hyperref}
\usepackage{xcolor}
\usepackage{enumitem}
\usepackage{titlesec}
\usepackage{parskip}
\usepackage{tcolorbox}
\usepackage{amsmath,amssymb}
\usepackage{booktabs}

\definecolor{accent}{HTML}{8B0000}
\definecolor{muted}{HTML}{555555}
\definecolor{lightbg}{HTML}{F7F4F0}

% Section: slide title
\titleformat{\section}{\large\bfseries\color{accent}}{\thesection.}{0.6em}{}
\titlespacing*{\section}{0pt}{1.2em}{0.2em}

% Subsection: structural labels (Talking points / If asked / etc.)
\titleformat{\subsection}{\normalsize\bfseries}{}{0em}{}
\titlespacing*{\subsection}{0pt}{0.6em}{0.1em}

\setlist[itemize]{leftmargin=1.2em,itemsep=0.15em,topsep=0.2em}

\newcommand{\timetag}[1]{\hfill\textcolor{muted}{\small[#1]}}
\newcommand{\highlight}[1]{\textcolor{accent}{\textbf{#1}}}

\title{\bfseries Milestone 2 --- Speaker Companion}
\author{15-minute presentation: per-slide guidance, background, and Q\&A backup}
\date{}

\begin{document}

\maketitle
\vspace{-2em}

\begin{tcolorbox}[colback=lightbg,colframe=accent!40,arc=2pt]
\textbf{Pacing.} 13 main-content slides over $\sim$15 minutes $\approx$ 45--75 seconds per slide. Spend most time on the highlighted headline-result and in-space placebo slides; move quickly through the methodology recap. Appendix slides are reserve material for Q\&A --- do not budget time for them. \textbf{Watch the clock at slide 11 (headline)}; if behind, compress Part 2 (slide 13) and go straight to thanks.
\end{tcolorbox}

% =====================================================================
\section{Title slide \timetag{10s}}

Open with the title. The framing is causal estimation of one specific event, not forecasting and not the network-topology project from the previous milestone. Mention that today is Part 1 of a two-part thesis.

% =====================================================================
\section{From the previous presentation \timetag{45s}}

\subsection{Talking points}
\begin{itemize}
  \item State the original framing: maritime trade network topology shaping oil prices, with a forecasting question and a propagation question.
  \item Pivot reason: feedback to focus on \highlight{causal estimation} rather than forecasting. Forecasting is hard and crowded; a credible counterfactual for one specific disruption is more tractable and policy-relevant.
  \item Reviewed several causal methods (event studies, structural VARs, DiD, network-shock propagation, SCM).
  \item SCM chosen for two reasons: \emph{inferential machinery} (in-space placebo, cross-event transfer) and \emph{interpretable result} (donor-weighted counterfactual on the observed price), and it produces an ATT directly rather than a shock decomposition.
\end{itemize}

\subsection{Behind the slide}
The maritime-network framing was scoped down rather than abandoned. The structural-VAR alternative (Kilian-style) was the most natural competitor but requires identification restrictions that are themselves contested and that have no standard precedent for chokepoint events. Event studies don't produce a continuous counterfactual; DiD needs a clear treated/control structure that doesn't exist for a single global commodity.

\subsection{If asked}
\begin{itemize}
  \item \textit{Why not just stick with the network framing?} Network propagation on oil prices is interesting but very hard to identify from observational data; the network-as-donor approach we considered ends up looking similar to SCM with engineered features. SCM gives a cleaner identification story now.
  \item \textit{Could the network framing come back in Part 2?} Yes --- network centrality or routing-disruption features could enter the donor pool, or layered as a propagation model on top of the ATT. Part 2 is open.
\end{itemize}

% =====================================================================
\section{Problem statement \timetag{60s}}

\subsection{Talking points}
\begin{itemize}
  \item \highlight{The question (highlighted in red box):} What was the Brent price impact of the 2026-02-01 Strait of Hormuz disruption?
  \item Estimand: ATT over a defined post-event window. Method: SCM on the observed price series.
  \item Part 2 is open; come back to it at the end.
  \item What is new: SCM-on-price design is rare (closest precedent is Becker, Pfeifer \& Schweikert 2021 on Austrian retail fuel); field default is structural VAR or scenario models.
  \item Method extensions vs the precedent: crude benchmark (not retail fuel), cross-asset donors (not cross-country), exogenous geopolitical treatment (not endogenous policy).
  \item \highlight{Two-event robustness design}: Russia 2022 used as a magnitude-validation case. SCM Russia counterfactual lands within \$3.98/bbl of EIA STEO's last pre-invasion structural forecast on the matched 151-day window ($\sim$5.7 pp on implied ATT).
\end{itemize}

\subsection{Behind the slide}
The two-event design is the most important methodological contribution. SCM on a global benchmark price is structurally different from SCM on a country-level outcome --- there's only one Brent. The Russia case validates that the design produces magnitudes consistent with an independent structural forecast, before applying it to the focal Hormuz case where no such external anchor exists (Hormuz is too recent).

\subsection{If asked}
\begin{itemize}
  \item \textit{Why is there no peer-reviewed SCM precedent on crude?} A combination of recency (price-impact SCM applications are mostly post-2020), data availability (good donor pools for global commodities are non-obvious), and convention (the oil literature defaults to structural VAR for identified shock decomposition; SCM is more common in policy-evaluation settings). We don't claim "no peer-reviewed work has tried it," just that the closest mechanical precedent is regulated retail fuel.
  \item \textit{What's the policy use?} SPR sizing, central-bank reaction function (oil shocks pass through to headline inflation), fiscal buffers, sovereign hedging.
\end{itemize}

% =====================================================================
\section{Brent spot price, post-COVID --- curated event timeline \timetag{45s}}

\subsection{Talking points}
\begin{itemize}
  \item Visual orientation: where Hormuz 2026 sits in the broader post-COVID Brent path. Roughly $\$60$/bbl in 2020, runs up through 2021--22 (Russia), declines through 2023--24, then the Hormuz spike in early 2026.
  \item Events colour-coded by category (demand, supply, policy, geopolitics) --- shows the data is event-rich, which is why donor cleanliness matters.
  \item Highlight the two focal events visually: \highlight{Hormuz 2026} (focal) and \highlight{Russia 2022} (calibration).
\end{itemize}

\subsection{If asked}
\begin{itemize}
  \item \textit{Why these two events specifically?} Both are large, well-measured, recent geopolitical disruptions. Hormuz 2026 is the focal event of interest; Russia 2022 is the most recent large oil shock with a long enough post-window to be well-measured, used as a magnitude-validation case.
  \item \textit{Why not Abqaiq 2019 or others?} Discussed in donor\_catalog.md: pre-windows for older events overlap with Hormuz's, breaking the cross-event independence we exploit in $\S$5f. Russia 2022 is independent of the Hormuz pre-window.
\end{itemize}

% =====================================================================
\section{Methodology --- ensemble of 5 models \timetag{60s}}

\subsection{Talking points}
\begin{itemize}
  \item Ensemble of 5 models spanning the SCM design space:
  \begin{itemize}
    \item \textbf{Convex SCM} (Abadie, Diamond \& Hainmueller 2010): canonical, no extrapolation.
    \item \textbf{Augmented SCM} (Ben-Michael, Feller \& Rothstein 2021): ridge bias correction to allow extrapolation when convex hull is too tight.
    \item \textbf{Elastic-net} (Doudchenko \& Imbens 2016): $L_1+L_2$ regression, sparse signed weights.
    \item \textbf{XGBoost} (Chen \& Guestrin 2016): nonlinear, captures interactions.
    \item \textbf{Bayesian Ridge}: Bayesian linear regression with normal prior. Note: this is NOT the full state-space BSTS of Brodersen 2015, just the regression component (clarified in methodology.md).
  \end{itemize}
  \item \highlight{Why an ensemble:} the headline is the median across models; IQR is the model-uncertainty band. No single model assumption is privileged.
\end{itemize}

\subsection{Behind the slide}
The ensemble spans key SCM regularization assumptions: simplex (convex SCM), ridge (ASCM, BR), L1+L2 (Elastic-net), tree-based (XGBoost). Disagreement across models is itself informative: tight agreement = factor structure is well-identified; wide spread = model-uncertainty band that the headline IQR captures honestly.

\subsection{Model formulas --- what each one fits}

Let $y_t$ = log-Brent at time $t$ and $y_{j,t}$ = log-price of donor $j$. All models produce a synthetic counterfactual $\hat{y}_t$; the post-event gap is $100 \cdot (\exp(\bar{y}_{post} - \bar{\hat{y}}_{post}) - 1)$ in percent.

\begin{itemize}
  \item \textbf{Convex SCM.} Synthetic is a \emph{simplex-constrained weighted average} of donors:
    \[
      \hat{y}_t = \sum_{j=1}^{J} w_j\, y_{j,t}, \qquad w_j \geq 0, \quad \sum_j w_j = 1.
    \]
    Read like portfolio shares: ``synthetic Brent is $w_A$ of donor A $+ w_B$ of donor B $+ \dots$''. Cannot go outside the donor envelope. Sparse: typically $\sim$5 donors get non-trivial weight, the rest pinned to zero.
    \emph{Feature selection:} the simplex constraint behaves like an $L_1$ ball --- the SLSQP optimizer drives most $w_j$ to exactly zero, so ``best features'' = the small set of donors whose convex combination best reproduces Brent in the pre-period (lowest pre-RMSPE). No external selection step; sparsity is mechanical.
  \item \textbf{Augmented SCM.} Convex SCM constrains the donor weights to be non-negative and sum to 1 --- interpretable but rigid. ASCM keeps those convex weights and adds a second ridge-regression layer that \emph{extrapolates the donors' contribution beyond the simplex} --- letting the synthetic borrow more (or less) from each donor than Convex SCM allowed:
    \[
      \hat{y}_t = \underbrace{\sum_j w_j\, y_{j,t}}_{\text{convex base}} + \underbrace{\sum_j \eta_j\, y_{j,t}}_{\text{ridge extrapolation}}, \qquad \eta = (X^\top X + \lambda I)^{-1} X^\top r,
    \]
    where $r$ is the pre-period gap left over by convex SCM and $X$ is the donor matrix. The effective weight per donor is $w_j + \eta_j$ --- the $\eta_j$ are unconstrained (can be negative or exceed 1), which is what produces the extrapolation outside the donor hull. The hyperparameter $\lambda$ = \texttt{ridge\_lambda} controls how much extrapolation is permitted: large $\lambda$ shrinks the ridge layer toward zero and recovers Convex SCM; small $\lambda$ gives the ridge layer free rein.
    \emph{Feature selection:} sparsity is inherited from the convex base ($w_j$) only; the ridge layer $\eta_j$ touches every donor (no selection). So ``best features'' = same as Convex SCM, plus every donor contributing a small ridge correction.
  \item \textbf{Elastic-net.} Linear regression with combined $L_1 + L_2$ penalty:
    \[
      \hat{y}_t = \beta_0 + \sum_j \beta_j\, y_{j,t}, \qquad
      \min_\beta \;\tfrac{1}{2N}\|y - X\beta\|_2^2 \;+\; \alpha \Big[\rho \|\beta\|_1 + \tfrac{1-\rho}{2}\|\beta\|_2^2\Big].
    \]
    Coefficients $\beta_j$ are signed and unconstrained; the $L_1$ term drives some to exactly zero (sparsity), the $L_2$ term shrinks the rest. $\alpha$ = overall penalty strength; $\rho$ = \texttt{l1\_ratio} (0 = pure ridge, 1 = pure lasso, 0.5 = balanced).
    \emph{Feature selection:} explicit via the $L_1$ term --- donors with weak partial correlation against Brent's residual get $\beta_j = 0$. Higher $\rho$ $\to$ sparser selection. ``Best features'' = donors whose marginal contribution survives the L1 threshold after accounting for collinearity with already-selected donors.
  \item \textbf{XGBoost.} Sum of $K$ shallow regression trees, fit sequentially to residuals:
    \[
      \hat{y}_t = \sum_{k=1}^{K} f_k(\mathbf{y}_{\cdot, t}), \qquad f_k \in \{\text{trees of depth} \leq d\},
    \]
    where each tree $f_k$ minimizes the residual after the previous $k-1$ trees, with $L_2$ regularization $\lambda$ on leaf weights and step size $\eta$ (learning rate). \emph{Not} a closed-form regression: a tree at each step picks the donor and threshold that best reduces the residual. Captures nonlinear interactions but cannot extrapolate outside the training range of any donor.
    \emph{Feature selection:} greedy and built-in. At every split, the tree-builder scores every donor's best threshold by gain reduction and picks the maximum; donors that never get picked have zero importance. ``Best features'' = donors with highest cumulative split-gain across all $K$ trees (reported as XGBoost's \texttt{feature\_importances\_}).
  \item \textbf{Bayesian Ridge.} Same regression equation as OLS, but with a Gaussian prior placed on the coefficients $\beta$ (and Gamma priors on the precisions):
    \[
      \underbrace{y_t = \beta_0 + \sum_j \beta_j\, y_{j,t} + \varepsilon_t, \quad \varepsilon_t \sim \mathcal{N}(0, 1/\alpha)}_{\text{data model (identical to OLS)}}
    \]
    \[
      \underbrace{\beta_j \mid \lambda \sim \mathcal{N}(0,\, 1/\lambda), \qquad \lambda \sim \text{Gamma}(\cdot), \qquad \alpha \sim \text{Gamma}(\cdot)}_{\text{prior on the unknowns (what makes it Bayesian)}}
    \]
    The features $y_{j,t}$ enter the same way they would in OLS --- as known data inside the regression sum. The Bayesian part is purely the prior on the unknowns ($\beta$, $\alpha$, $\lambda$). Posterior over $\beta$ is closed-form (conjugate), so $\hat{y}_t$ comes with a native posterior standard deviation, no bootstrap needed. $\lambda$ (the weight-precision prior) is the regularization knob --- larger $\lambda$ $\to$ tighter prior around zero $\to$ stronger shrinkage; $\alpha$ is the noise-precision prior. MAP-equivalent to classical Ridge with penalty $\lambda/\alpha$.
    \emph{Feature selection:} \textbf{none --- all donors retained}. The Gaussian prior shrinks coefficients smoothly toward zero but never sets any to exactly zero, unlike $L_1$ models. ``Best features'' is approximated by the PIP proxy $|\hat\beta_j| / \sigma_{\hat\beta_j}$ (the posterior $z$-score), which ranks donors by how confidently their coefficient differs from zero. To get \emph{actual} variable selection in a Bayesian framework, you would need spike-and-slab priors (the full BSTS of Brodersen 2015), which this implementation does not include.
\end{itemize}

\subsection{Tuned hyperparameters --- chosen by walk-forward CV val RMSE}

Grids tightened per Cawley \& Talbot (2010) to limit winner's-curse on val RMSE. Values below are the \emph{chosen} settings per event after the 5-fold walk-forward search.

\begin{itemize}
  \item \textbf{Convex SCM} --- no tuned hparams. Fixed: 120 random $V$ matrices drawn from a Dirichlet, weights re-fit per $V$, best by pre-period RMSPE retained. Same for both events.
  \item \textbf{Augmented SCM} (grid: $\lambda \in \{0.01, 0.1, 1, 10, 100\}$, 5 candidates):
  \begin{itemize}
    \item Hormuz: \texttt{ridge\_lambda = 10.0} (heavier extrapolation control).
    \item Russia:  \texttt{ridge\_lambda = 0.1} (light extrapolation control --- the convex residual is bigger so ridge needs more freedom).
  \end{itemize}
  \item \textbf{Elastic-net} (grid: $\alpha \in \{0.001, 0.01, 0.05, 0.1, 0.5\} \times \rho \in \{0.2, 0.5, 0.8\}$, 15 candidates):
  \begin{itemize}
    \item Hormuz: \texttt{alpha = 0.001}, \texttt{l1\_ratio = 0.8} (sparser, lasso-leaning).
    \item Russia:  \texttt{alpha = 0.001}, \texttt{l1\_ratio = 0.2} (denser, ridge-leaning).
  \end{itemize}
  \item \textbf{XGBoost} (grid: \texttt{max\_depth} $\in \{2,3,4\}$ $\times$ \texttt{learning\_rate} $\in \{0.001, 0.005, 0.01, 0.03, 0.1\}$ $\times$ \texttt{reg\_lambda} $\in \{1, 10\}$, 30 candidates; \texttt{n\_estimators=1000} with early stopping after 20 rounds):
  \begin{itemize}
    \item Hormuz: \texttt{max\_depth=2}, \texttt{learning\_rate=0.1}, \texttt{reg\_lambda=10.0}.
    \item Russia:  \texttt{max\_depth=2}, \texttt{learning\_rate=0.1}, \texttt{reg\_lambda=1.0}.
  \end{itemize}
  Both events stay at depth 2 --- only second-order interactions, consistent with the small-data regime.
  \item \textbf{Bayesian Ridge} (grid: $\lambda \in \{10^{-8}, 10^{-6}, 10^{-4}, 10^{-2}, 1\} \times \alpha \in \{10^{-8}, 10^{-6}, 10^{-4}\}$, 15 candidates):
  \begin{itemize}
    \item Hormuz: $\lambda_1=\lambda_2=1.0$, $\alpha_1=\alpha_2=10^{-8}$ (strong weight shrinkage, near-flat noise prior).
    \item Russia:  $\lambda_1=\lambda_2=1.0$, $\alpha_1=\alpha_2=10^{-8}$ (same).
  \end{itemize}
  Both events converge on the strongest prior in the grid --- the donor regression is high-variance at $N \approx 420$ and shrinkage helps both.
\end{itemize}

\subsection{If asked}
\begin{itemize}
  \item \textit{Why not include BSTS proper?} The implementation in lib/models.py is \texttt{sklearn.linear\_model.BayesianRidge} --- Bayesian linear regression on donor levels, not the full state-space BSTS with trend + seasonal + spike-and-slab. Documented honestly in methodology.md $\S$4 ``Bayesian Ridge vs full BSTS.''
  \item \textit{Why not include simpler models like OLS?} OLS without regularization overfits at $N \approx 420$ obs and 21 donors. Elastic-net at low $\alpha$ approximates OLS while preventing overfit.
  \item \textit{Why does ASCM pick very different $\lambda$ across the two events?} Russia has a larger convex-residual to absorb (Brent's pre-period $RMSPE \approx 0.106$ vs Hormuz's $\sim$0.066), so the ridge fix-up needs more freedom (lower $\lambda$) to fit it. Hormuz's convex fit is tight enough that a heavy $\lambda$ doesn't hurt and protects against extrapolation.
  \item \textit{Why does Elastic-net pick opposite \texttt{l1\_ratio} across events?} Hormuz's signal is concentrated in a few donors (sparser fit preferred $\to$ $\rho{=}0.8$); Russia's signal is spread across many correlated donors (ridge-leaning $\to$ $\rho{=}0.2$).
\end{itemize}

% =====================================================================
\section{Model training and validation \timetag{75s}}

\subsection{Talking points}
\begin{itemize}
  \item Pre-event windows: Hormuz 2024-06 $\to$ 2026-01 (423 obs); Russia 2020-07 $\to$ 2022-02 (421 obs). Post-event: $\sim$59 obs Hormuz / 151 obs Russia.
  \item \highlight{Training:} 5-fold expanding-window walk-forward CV (Hyndman \& Athanasopoulos 2018; Bergmeir \& Benítez 2012). Hyperparameters chosen by pooled val RMSE across folds; final fit refits on full pre-event window.
  \item Validation battery: model fit (walk-forward), SUTVA (donor audit), inference (in-space placebo + cross-event transfer).
\end{itemize}

\subsection{Behind the slide}
The walk-forward CV is a recent refactor (previously was a single 80/20 hold-out, which is a time-series train-test split, not true walk-forward CV per Hyndman). The 5-fold expanding-window scheme is more robust because it averages over 5 forecast windows rather than concentrating all val mass on the last 20\%. Headline numbers shifted modestly under the refactor: Hormuz unchanged at 43\%; Russia ensemble median 25.1\% $\to$ 21.7\%.

\subsection{If asked}
\begin{itemize}
  \item \textit{Why expanding-window rather than sliding-window or pure k-fold?} K-fold breaks time-series structure (uses future to predict past). Sliding-window discards early data. Expanding-window matches the operational reality (you have more data as you go).
  \item \textit{Hyperparameter selection bias?} Yes: hparams are picked to minimize val RMSE, so reported val RMSE is biased downward. Documented in validation.md $\S$5a as a ``winner's-curse'' caveat. Mitigation: tightened grids (each model has at most 9 candidates) per Cawley \& Talbot 2010.
  \item \textit{Why not a separate test set?} Pre-event window is sample-size-constrained. Splitting off a test set would shrink training meaningfully. Formal inference is delegated to the in-space placebo (which is genuinely out-of-sample at the unit level).
\end{itemize}

% =====================================================================
\section{Donor selection \timetag{75s}}

\subsection{Talking points}
\begin{itemize}
  \item \highlight{21 audited donors out of 33 candidates.} The full 33-donor catalog was audited for both pre-treatment co-movement with Brent AND treatment cleanliness (SUTVA). 12 dropped due to event-specific exposure.
  \item Categories: metals (3), agriculturals (4), equities (3), FX (8), rates/credit (2), volatility (1).
  \item Inclusion rule: pre-treatment co-movement through a \emph{common global factor} AND treatment cleanliness (no Russia/Ukraine or Persian Gulf exposure).
  \item Same pool used for both events, enabling the cross-event weight transfer test ($\S$5f).
\end{itemize}

\subsection{Behind the slide}
Donor selection is the SCM-canonical SUTVA defense (Abadie 2021 $\S$3.1: substantive grounds, not statistical tests). The audit reasoning is in donor\_catalog.md. The dropped 12 includes: Russia commodities (Wheat, Palladium, Russian-channel FX like EUR), risk-off proxies that loaded on Russia specifically (DXY, BTC due to sanctions-evasion narrative), and others.

\subsection{If asked}
\begin{itemize}
  \item \textit{What's the full list of dropped donors?} See backup slide ``Excluded donors.'' 12 Russia-exposed (H or M tier): Copper, IronOre, Palladium, Soybeans, Wheat, Corn, EM\_Eq, EUR, GBP, DXY, US10Y, BTC. Plus another set of categorical exclusions (energy complex itself: WTI, NatGas, refined products; oil-currency: NOK, CAD; etc.) that never entered the candidate pool.
  \item \textit{What if a clean donor is actually treated?} The model-agnostic battery in 01.5\_Donor\_Cleanliness checks this empirically. Hormuz: 32/33 agreement between audit and tests; Russia: 20/33. Discrepancies are concurrent macro events (CNY's China zero-COVID), not Russia-treatment.
  \item \textit{Why 21 specifically and not 24 or 18?} The shared pool is the intersection of Russia-clean (21 of 33) and Hormuz-clean (33 of 33). Russia is the binding constraint.
\end{itemize}

% =====================================================================
\section{Donor co-movement with Brent --- audited pool \timetag{30s}}

\subsection{Talking points}
\begin{itemize}
  \item Visual: 21 audited donors rebased to 100 at pre-window start, with Brent in black on top. Dashed line marks $T_0$.
  \item Show the audited donors are visibly correlated with Brent in the pre-event window --- the SCM has something real to fit.
  \item Move quickly; this is an orientation slide, not an analytical one.
\end{itemize}

% =====================================================================
\section{Validation results --- model fit \timetag{60s}}

\subsection{Talking points}
\begin{itemize}
  \item Ensemble paths (left Hormuz, right Russia): all 5 model synthetics overlay observed Brent in the pre-event window, then fan out post-event. Visual evidence of pre-event fit AND post-event divergence.
  \item Walk-forward CV table: \highlight{train RMSE / val RMSE / val-train ratio per model, both events}. Numbers are in \textbf{log-Brent units} (the fit target is log-price levels, not log-returns).
  \item Convex SCM val/train ratio is 1.26 Hormuz / 1.16 Russia (clean). On Russia, ASCM 2.04 (just over 2$\times$ heuristic threshold), XGBoost 3.00, and Bayesian Ridge 3.07 are the overfit-flagged cases.
  \item Hormuz: 4 of 5 models clean; only Bayesian Ridge flagged.
  \item Ensemble is robust to individual model overfit because median across 5 models smooths.
\end{itemize}

\subsection{Hormuz vs Russia --- the headline distinction}

\textbf{Hormuz is the focal causal estimate.} Clean post-period with no concurrent macro confounder; all 5 models converge in the 38--50\% range; in-space placebo hits the floor for every cell. Identification holds.

\textbf{Russia is the methodology-validation case.} Identification is partially weakened because the post-window (Feb--Sep 2022) overlaps a major macro regime shift:
\begin{itemize}\setlength\itemsep{0.05em}
  \item Fed liftoff March 2022 (+25 bp), then +75 bp $\times$ 3 consecutive hikes.
  \item DXY surged $\sim$95 $\to$ $\sim$115.
  \item JPY collapsed $\sim$115 $\to$ $\sim$145 (BoJ stayed on YCC, carry trade unwound).
  \item Bond crash (TLT $-30\%$), global risk-off, credit-spread widening.
\end{itemize}
These hit the donor pool (JPY, KRW, MXN, TLT, HYG, WorldEq, $\dots$) even though the qualitative audit certified them ``Russia-clean.'' They moved for \emph{macro} reasons, not Russia-spillover.

\textbf{The cleanest read of the Russia premium is in the first 2--3 weeks post-invasion}: Brent went \$95 (Feb 23) $\to$ \$130+ (Mar 8), well before Fed hiking became dominant. A shorter post-window would isolate the Russia effect more cleanly --- one-line config change in \texttt{lib/config.py} if a reviewer asks.

\textbf{Defense for the Russia case in spite of weakened identification:} (a) EIA STEO agreement within \$3.98/bbl on the matched window, (b) sparse-model convergence around 22--30\% (which exclude the macro-confounded donors via $L_1$/simplex sparsity), (c) well-documented historical magnitude. The Russia ATT is \emph{not} a clean point estimate; it's a credibility check that the methodology recovers a believable magnitude on a known disruption.

\subsection{Why Bayesian Ridge's Russia synthetic ``always increases''}

The visible failure mode: BR's Russia synthetic tracks Brent's post-event rally rather than staying flat, eating the Russia premium and giving a degenerate $+0.5\%$ ATT. Three stacked reasons:
\begin{itemize}\setlength\itemsep{0.05em}
  \item \textbf{Dense regression, no feature selection.} BR retains all 21 donors with non-zero coefficients (unlike Convex SCM / Elastic-net, which zero out most). When the macro-confounded donors moved up post-event, BR's synthetic moved up with them.
  \item \textbf{Russia's macro confounder pushes those donors in the same direction as Brent.} Fed hiking $\to$ dollar surge $\to$ JPY collapse, bond crash, equity drawdown. Dense weights pick all that up; sparse weights drop most of it.
  \item \textbf{BR is overfit on Russia} (val/train ratio 3.07). Overfit models amplify input-feature movements at test time.
\end{itemize}
Hormuz post-period has no comparable macro shock, so BR's Hormuz synthetic stays flat alongside the other models and the ATT is in line ($+43.6\%$).

\subsection{Interpreting the RMSE numbers --- log-Brent to dollars}

The fit target is log-Brent (not log-returns), so RMSE is in log-price units. Conversion: a residual $\sigma$ in log-units $\to$ multiplicative error $e^\sigma$; for small $\sigma$, $e^\sigma - 1 \approx \sigma$, so \textbf{log-RMSE $\approx$ fractional error to first order}.

\begin{center}
\begin{tabular}{@{}c c c c c@{}}
  \toprule
  Log RMSE & $e^\sigma - 1$ & \% typical err. & at Brent \$70/bbl & at Brent \$100/bbl \\
  \midrule
  0.031 & 3.15\%  & $\sim$3\%  & $\pm$\$2.20/bbl  & $\pm$\$3.10/bbl \\
  0.059 & 6.08\%  & $\sim$6\%  & $\pm$\$4.25/bbl  & $\pm$\$6.10/bbl \\
  0.075 & 7.79\%  & $\sim$8\%  & $\pm$\$5.45/bbl  & $\pm$\$7.80/bbl \\
  0.109 & 11.52\% & $\sim$12\% & $\pm$\$8.10/bbl  & $\pm$\$11.50/bbl \\
  0.127 & 13.55\% & $\sim$14\% & $\pm$\$9.50/bbl  & $\pm$\$13.50/bbl \\
  \bottomrule
\end{tabular}
\end{center}

Reading the slide table with this:
\begin{itemize}\setlength\itemsep{0.05em}
  \item Convex SCM Hormuz train 0.059 $\to$ $\sim$6\% typical pre-period error ($\sim$\$4--5/bbl at Hormuz prices).
  \item Convex SCM Russia val 0.127 $\to$ $\sim$14\% val error ($\sim$\$9--11/bbl at Russia prices). The widest fit gap on the board, reflecting that the simplex constraint cannot span Brent's 2020--22 range with the audited donor set.
  \item Bayesian Ridge train 0.031 / val 0.096 $\to$ $\sim$3\% in training but $\sim$10\% in val. The gap between training-period precision and val-period error is the overfit signature.
\end{itemize}

\subsection{Behind the slide}
The synthetic represents ``what Brent would have done given how the donors moved.'' If donors moved a lot (Russia post-period) the synthetic naturally moves too; sparse models drop the macro-driven donors and the synthetic stays flatter. BR keeps them in and the synthetic rises with Brent. The model spread on Russia (0.5\% to 36.8\%) is itself diagnostic --- it tells you the macro confounder is generating model-specific identification problems, not just numerical noise.

\subsection{If asked}
\begin{itemize}
  \item \textit{Why don't you exclude the overfit models?} The heuristic threshold (val/train $>2$) is just a flag, not a gate. Auto-excluding would impose a literature-unsupported cutoff and would tighten the IQR artificially. Defense is the IQR-as-uncertainty-band convention.
  \item \textit{Is the validation downward biased by hyperparameter selection?} Yes; documented. Magnitude: the bias is small ($\sim$10\% on val RMSE) because grids are tight per Cawley \& Talbot.
  \item \textit{Doesn't the 2022 macro confounder invalidate the Russia SCM?} Weakens, doesn't fully invalidate. Identification claim becomes ``Russia premium + asymmetric macro impact'' rather than pure Russia ATT. Defended by (a) EIA STEO agreement within \$3.98/bbl, (b) convergence of sparse models 22--30\%, (c) historical record \$95 $\to$ \$130, (d) Hormuz works cleanly as the methodology validation. We use Russia as a magnitude-credibility check, not as a clean causal ATT.
  \item \textit{Why is RMSE in log-Brent and not absolute Brent?} Standard SCM practice (Abadie 2010 fits log-levels). Log scale makes the relative error comparable across price regimes (a \$5/bbl miss at \$30 is much worse than at \$120). Also makes the SCM gap interpretable as a percentage directly.
  \item \textit{Why log-levels and not log-returns?} SCM is about counterfactual \emph{paths} (where would Brent's level be), not about return prediction. Fitting returns and then integrating compounds errors; fitting levels directly produces the path you want to report.
\end{itemize}

% =====================================================================
\section{Validation results --- in-space placebo \timetag{90s}}

\subsection{Talking points}
\begin{itemize}
  \item \highlight{Plain-English summary:} re-run the same SCM machinery 21 more times, each time pretending a different donor was the treated unit instead of Brent. If Brent's gap is bigger than every fake gap the procedure invents for these untreated donors, the effect is unlikely to be something the model would have produced by accident.
  \item Abadie's randomization inference: refit treating each donor as the placebo treated unit, compute Brent's rank in the post/pre RMSPE distribution.
  \item \highlight{Permutation floor}: 1/22 $\approx 0.045$ for 21-donor pool; 1/34 $\approx 0.029$ for 33-donor full pool.
  \item Hormuz: \highlight{all 5 models hit the floor under both pools} --- Brent ranks 1 of 22 in each. Strongest possible permutation signal.
  \item Russia: not decisively rejected. Only XGBoost reaches floor at 21-donor (0.045) and slips to 0.147 at 33-donor. Other models p $\sim$ 0.4--0.7.
\end{itemize}

\subsection{Behind the slide}
The two-pool design (21-donor and 33-donor full) is a robustness check. The 21-donor floor is too coarse to distinguish ``Brent is rank 1'' from anything below 0.045. The 33-donor pool with audit-flagged donors mixed back in tests whether the rank-1 finding holds when ``noisier'' donors expand the placebo distribution. Hormuz is robust; Russia is not.

The Russia weakness is the known limitation of in-space placebo at pre-periods with high donor volatility (Russia's pre-period 2020-07 $\to$ 2022-02 spans COVID recovery + 2021 reflation rally, both of which moved donors). This widens the placebo distribution and prevents Brent's post-period ratio from sitting in the tail.

\subsection{If asked}
\begin{itemize}
  \item \textit{Is XGBoost's Russia rejection meaningful?} It clears at 21-donor floor but fails at 33-donor (p $\sim$ 0.147), so the rejection is fragile. Russia's case for a real effect rests on (a) the historical record (well-documented \$95 $\to$ \$130 move), (b) the audit-supported SUTVA defense, and (c) the cross-event weight transfer (slide A5), not on the in-space placebo alone.
  \item \textit{Why does Hormuz have such tight p-values?} Brent's post/pre RMSPE ratios for Hormuz are 6 to 12 across models --- much bigger than any donor's. The donor pool is also more stable in the Hormuz pre-window (no COVID-style regime shifts), so the placebo distribution is tight.
  \item \textit{Permutation test theory: how do we know p-values are valid?} The randomization-inference framework (Fisher 1935; Lehmann \& Romano 2005) is what justifies SCM's p-values. Abadie 2010 $\S$4 explicitly grounds the placebo inference in this framework. Exact at the unit level (we enumerate all 22 or 34 placebo refits).
\end{itemize}

% =====================================================================
\section{External validation --- EIA STEO pre-invasion forecasts (Russia only) \timetag{75s}}

\subsection{Talking points}
\begin{itemize}
  \item Two-method cross-check: SCM Russia counterfactual vs EIA STEO's pre-invasion forecast of 2022 Brent, measured on the \emph{exact same 151-day treatment window} (Feb 24 $\to$ Sep 30 2022).
  \item Two STEO vintages: \highlight{Jan-22} (issued $\sim$Jan 11, 2022 --- 6 weeks pre-invasion) and \highlight{Feb-22} (issued $\sim$Feb 8, 2022 --- 16 days pre-invasion, last pre-invasion STEO).
  \item Result table: SCM ensemble-median counterfactual \$88.92/bbl. STEO Feb-22 step-function \$84.94/bbl. \highlight{Disagreement: only \$3.98/bbl, or 5.7 pp on implied ATT.}
  \item Three of five SCM models bracket STEO Feb-22 within $\pm$\$6.
  \item Implied ATTs: SCM-based 22.1\%, STEO-Feb-22-based 27.8\% ($\Delta \sim$ 5.7 pp), STEO-Jan-22-based 43.4\% (further in time = larger gap).
\end{itemize}

\subsection{Behind the slide}
SCM and STEO share no information path: SCM uses cross-asset co-movement; STEO uses EIA's structural supply/demand model. The agreement is independent-method evidence the SCM isn't producing a fantasy counterfactual. The disagreement direction matters too: SCM is more \emph{conservative} (lower implied ATT) because donor co-movement absorbs macro spillover (inflation, demand recovery) that the STEO doesn't.

This is uniquely available for Russia: Hormuz has no equivalent pre-event structural forecast because Hormuz happened in early 2026 and the methodology was developed after the event.

\subsection{If asked}
\begin{itemize}
  \item \textit{Why use STEO and not other forecasts?} STEO is the most prominent monthly oil forecast and publicly available with archive vintages. EIA Short-Term Energy Outlook archives at \texttt{eia.gov/outlooks/steo/archives}.
  \item \textit{Is STEO the ground truth?} No; STEO is another model. The agreement says two independent methods produce the same counterfactual, not that either is correct.
  \item \textit{What if STEO is just wrong?} Then the SCM might also be wrong in the same direction. The two-method agreement is necessary but not sufficient evidence. The qualitative case rests on the broader battery (placebo + cross-event transfer + donor audit).
\end{itemize}

% =====================================================================
\section{Headline result --- ATT estimate \timetag{90s}}

\subsection{Talking points}
\begin{itemize}
  \item \highlight{Headline numbers:}
  \begin{itemize}
    \item Hormuz 2026 ensemble median: \highlight{$\sim$43--44\%} chokepoint premium over the 3-month post-window. IQR [40.3, 44.2].
    \item Russia 2022 ensemble median: \highlight{$\sim$21--22\%} drift-adjusted Brent premium over the 7-month post-window. IQR [13.1, 23.8].
  \end{itemize}
  \item Hormuz IQR is tighter than Russia (4 pp vs 11 pp) --- donors are more stable in the Hormuz pre-period.
  \item Drift adjustment: subtract the pre-period slope $\times$ post-window length per Roth 2022 (pretest-with-caution framework).
  \item Hormuz: validated by in-space placebo at the floor under both pools across all 5 models, plus cross-event weight transfer (appendix). Russia: matches EIA STEO to within \$3.98/bbl.
\end{itemize}

\subsection{Behind the slide}
This is the punchline of Part 1. The Hormuz $\sim$43\% is the substantive headline; Russia $\sim$22\% is the methodology-validation case. The cross-event weight transfer (appendix) shows that the same factor structure that fits Russia transfers to Hormuz for the linear/sparse models (Convex SCM, ASCM, Elastic-net), giving regime-stability evidence for the focal estimate.

If asked for a single number for Hormuz: 43\%. If asked for a range: 38--50\% across models, 40--44\% interquartile.

\subsection{If asked}
\begin{itemize}
  \item \textit{How does 43\% compare to scenario-model estimates?} Dallas Fed WP 2609, Kiel Policy Brief 206, and IEA OMR April--May 2026 report scenario-model estimates for partial vs full Hormuz closure. The SCM is in the same order of magnitude as the partial-closure scenarios.
  \item \textit{What's the standard error?} The IQR is the model-uncertainty band; no formal SE is reported because SCM doesn't naturally have one. Inference is via the in-space placebo p-value.
  \item \textit{Is the Hormuz post-window long enough?} 59 obs is short by SCM convention. The cross-event weight transfer substitutes inferential depth.
\end{itemize}

% =====================================================================
\section{Part 2 --- ideas for the second half \timetag{45s}}

\subsection{Talking points}
\begin{itemize}
  \item Currently a placeholder for discussion. Possible directions (per the muted footnote):
  \begin{itemize}
    \item Pass-through of the Brent ATT to retail fuel prices and headline CPI.
    \item Direct comparison: SCM ATT vs scenario-model price paths (Dallas Fed, OIES, Kiel) on the same window.
    \item Severity-dependent extrapolation: the SCM gives one ATT for the realized disruption; what about counterfactual closure intensities?
    \item Alternative donor pool designs or maritime-network features re-entering as donor candidates.
  \end{itemize}
  \item \highlight{Open to suggestions from the audience.}
\end{itemize}

\subsection{Behind the slide}
The Part 2 slide is deliberately under-specified. The intent is to elicit feedback from the committee on direction. Have one or two preferred directions ready in case nobody pushes back: the most natural is the SCM vs scenario-model comparison (slide reuses the STEO machinery from Part 1).

% =====================================================================
\section{Thank you / Q\&A}

Open the floor. The appendix slides are reserve material; navigate to them on demand:

\begin{itemize}
  \item \textbf{Appendix A1 --- excluded donors:} 12 Russia-treated drops with audit reasons; categorical exclusions (energy complex, oil-currencies).
  \item \textbf{Appendix A2 --- donor importance:} per-model weight heatmaps (note: weights are not identical across events, see A5 nuance).
  \item \textbf{Appendix A3 --- donor SUTVA cleanliness:} permutation mean-shift test results (Chow 1960 hypothesis + Lehmann \& Romano 2005 permutation + Phipson \& Smyth 2010 correction). Hormuz 1/33 flagged (IronOre); Russia 1/33 (CNY). Audit agreement: Hormuz 32/33, Russia 20/33.
  \item \textbf{Appendix A4 --- in-time placebo and LOO:} combined diagnostics table. Tightest LOO: XGBoost ($\sim$5 pp both events). Widest: Bayesian Ridge Russia (24.9 pp, range crosses zero) --- the LOO-fragile case.
  \item \textbf{Appendix A5 --- cross-event weight transfer:} the $\S$5f test. Russia weights applied to Hormuz panel: Convex SCM transfers cleanly (+8 pp), ASCM and Elastic-net transfer with attenuation, XGBoost and Bayesian Ridge transfer pathologically (transferred pre-RMSPE 22$\times$ and 314$\times$ independent). The relevant subtlety: per-donor weights aren't stable across events (e.g., JPY moves from rank 12 to rank 1 in convex SCM), but the \emph{projection onto donor space} is. Cite Abadie \& L'Hour 2021 on convex SCM weight non-uniqueness.
\end{itemize}

% =====================================================================
\section*{General Q\&A backup material}

\subsection{Methodological choices to be ready to defend}
\begin{itemize}
  \item \textbf{Why 5 models and not 1?} Ensemble spans the SCM regularization design space; no single model assumption is privileged; IQR is the honest uncertainty band.
  \item \textbf{Why log-Brent and log-donors?} Stationarity (log-returns are approximately stationary; levels are not), interpretability (log differences $\approx$ percentage changes), variance stabilization.
  \item \textbf{Why a shared 21-donor pool for both events?} Enables the cross-event weight transfer test (the strongest single methodology generalization test). Alternative per-event pools would give somewhat different headline numbers but break the cross-event design.
  \item \textbf{Why Russia 2022 specifically as the validation case?} Most recent large oil shock with a long post-window; independent of Hormuz pre-period; well-documented; matches the Hormuz design template.
  \item \textbf{Why pre-windows starting 2020-07 (Russia) and 2024-06 (Hormuz)?} 18--20 month pre-windows balance sample size against regime stability. Earlier starts include too much COVID disruption (Russia) or pre-2024 macro instability (Hormuz). Documented in PRE\_WINDOWS at lib/config.py:29.
\end{itemize}

\subsection{Numbers to know cold}
\begin{itemize}
  \item Hormuz headline drift-adjusted ATT: \textbf{43\%} (IQR 40--44).
  \item Russia headline drift-adjusted ATT: \textbf{21--22\%} (IQR 13--24).
  \item Actual Russia Brent mean over Feb 24 $\to$ Sep 30 2022: \$108.54/bbl.
  \item SCM ensemble median Russia counterfactual: \$88.92/bbl.
  \item EIA STEO Feb-22 forecast for the same window: \$84.94/bbl.
  \item Hormuz pre-event window: 423 obs / Russia pre-event: 421 obs.
  \item Hormuz post-event obs: $\sim$59 (the binding constraint on in-event inference).
  \item In-space placebo floors: 1/22 = 0.045 (21-donor); 1/34 = 0.029 (33-donor).
\end{itemize}

\subsection{Hardest critiques}
\begin{itemize}
  \item \textbf{``Your Hormuz IQR is 4 pp --- too narrow given 60 obs.''} The IQR reflects model agreement, not statistical uncertainty. The structural uncertainty is in the in-space placebo (floor) and cross-event transfer (degraded RMSPE for nonlinear models). Honest reading: 38--50\% across all models, 40--44\% across the agreeable models.
  \item \textbf{``Cross-event weight transfer fails for 2 of 5 models.''} It fails for the dense / regime-dependent models (XGBoost, Bayesian Ridge) and works for the sparse models (Convex SCM, ASCM, Elastic-net). This is mechanically expected: sparse weights pin to factor-stable donors. The Hormuz headline rests on the sparse models, which transfer.
  \item \textbf{``In-space placebo for Russia is weak.''} Yes. Russia's defense is multi-layered: EIA STEO external validation, historical record (\$95 $\to$ \$130 move is documented), audit-supported SUTVA, cross-event weight transfer. The in-space placebo is one of four pieces of evidence.
  \item \textbf{``You changed the in-time placebo to use defaults instead of tuned hparams.''} Yes, intentionally: the fake-post period overlaps the val window used for hparam tuning, so tuned hparams would induce second-order leakage. Defaults are the literature-defensible choice for an unselection-bias test. Documented in validation.md $\S$5e (ii).
\end{itemize}

\vspace{2em}
\hrule
\vspace{0.5em}
\textcolor{muted}{\small\itshape Companion compiled from docs/methodology.md, docs/validation.md, docs/results.md, docs/donor\_catalog.md, and docs/motivation.md as of commit \texttt{8721f35}. For full details on any item, the source-of-truth documents are in \texttt{docs/}.}

\end{document}
